In this chapter, we present \textbf{CoMesh}~\cite{comesh}, the first fully decentralized, hub-less control plane for Sense-Trigger-Actuate (STA) operations in large commercial edge IoT mesh networks. Rather than routing coordination through a central hub or cloud, CoMesh keeps all decision-making local to the mesh by organizing smart devices into overlapping subgroups called \emph{coteries} that collectively monitor devices and manage routines via zero-message exchange protocols, quorum-based agreement, epoch-based migration, and locality-sensitive hashing (LSH). The result is a self-healing architecture that adapts autonomously to device failures, churn, and mobility without any human intervention or centralized oversight, realizing a vision of edge-first smart environments that remain fully operational without external infrastructure.

This chapter is organized as follows. \Cref{comesh:section:intro} introduces the problem and surveys the commercial edge landscape; \Cref{comesh:sec:motivation} presents CoMesh's design overview. \Cref{comesh:sec:system-model} defines the system model. \Cref{comesh:section:design} describes CoMesh's coterie mechanism, including zero-message member selection, quorum-based agreement, epoch-based migration, and LSH-based locality. \Cref{comesh:section:mngmt} covers device and routine management, including conflict-free actuation via distributed locks. \Cref{comesh:sec:analysis} provides formal proofs of CoMesh's safety, liveness, and progress properties. \Cref{comesh:sec:expts} reports our Raspberry Pi deployment results and \Cref{comesh:sec:simres} presents large-scale simulation results using real building maps. \Cref{comesh:section:related_work} discusses related work and \Cref{comesh:section:conclusion} concludes.
